% SPDX-License-Identifier: GPL-3.0-or-later
% Copyright (C) 2026 Alessandro Greco
\chapter{Криптография, необходимая для Tor}\label{cap:05-crittografia}
\aggiornatoal{23 сентября 2026 г.}{tor 0.4.9.13; коммит torspec \texttt{8e1f245}}

\begin{obiettivi}
К концу этой главы вы сможете:
\begin{itemize}[nosep]
  \item определять криптографические примитивы, используемые Tor, и назначение каждого из них;
  \item различать шифрование, контроль целостности, выработку ключей, обмен ключами и цифровые подписи;
  \item понимать, почему Tor использует несколько семейств хеш-функций (SHA-1, SHA-2, SHA-3, Blake2b) и чем они различаются;
  \item уверенно читать формулы в последующих главах.
\end{itemize}
\end{obiettivi}

Эта глава --- набор инструментов, а не курс криптографии. Здесь с минимально необходимой теорией вводятся средства, которые встретятся нам в части~II, и каждый раз указывается, где именно они используются. В приложении~B приведены дополнительные сведения для читателей, желающих глубже изучить тему. Все примеры можно повторить в Python с помощью стандартного модуля \texttt{hashlib} и библиотеки \texttt{cryptography} --- тех же инструментов, которые используются в лабораторных работах.

% ---------------------------------------------------------------------------
\section{Пять задач}\label{sec:cripto-compiti}

В протоколе вроде Tor криптография решает пять разных задач:
\begin{description}
  \item[Конфиденциальность] Никто, кроме предполагаемых получателей, не должен иметь возможности читать данные: \emph{шифры}.
  \item[Целостность] Никто не должен иметь возможности незаметно изменить данные: \emph{хеш-функции} и \emph{коды аутентификации сообщений} (MAC).
  \item[Согласование ключей] Две стороны, изначально не имеющие общего секрета, должны установить его, обмениваясь сообщениями по открытому каналу: \emph{обмен ключами} Диффи--Хеллмана.
  \item[Аутентификация] Одна сторона должна доказать владение определённой идентичностью: \emph{цифровые подписи}.
  \item[Выработка] Из общего секрета необходимо получить несколько независимых ключей: \emph{функции выработки ключей} (KDF).
\end{description}
\begin{figure}[htbp]
\centering
% SPDX-License-Identifier: GPL-3.0-or-later
% Copyright (C) 2026 Alessandro Greco
%% Русское альтернативное описание EPUB.
\altfig{Пять криптографических задач Tor: конфиденциальность с AES-CTR для ячеек и AES-GCM либо ChaCha20-Poly1305 для каналов TLS (главы 10 и 7); целостность с дайджестами SHA-1 и SHA3-256 и тегами CGO на основе POLYVAL (глава 10); согласование ключей X25519 в ntor и ntor-v3 (глава 9); аутентификация Ed25519 и устаревающим RSA в сертификатах и консенсусе (главы 7 и 12); вывод ключей HKDF-SHA256 и SHAKE-256 (глава 9).}
% Пять криптографических задач Tor и используемые примитивы.
\begin{tikzpicture}[font=\sffamily\scriptsize,
  compito/.style={draw=torviola,fill=torviola!8,rounded corners=2pt,minimum width=2.6cm,minimum height=0.7cm},
  prim/.style={draw=torverde!60!black,fill=torverde!12,rounded corners=2pt,text width=3.9cm,align=center,minimum height=0.7cm},
  dove/.style={font=\sffamily\tiny,text=torgrigio,anchor=west}]
  \foreach \i/\c/\p/\d in {%
    0/Конфиденциальность/{AES-CTR; TLS с AES-GCM или ChaCha20-Poly1305}/{ячейки (гл.~10), каналы (гл.~7)},
    1/Целостность/{SHA-1 и SHA3-256, POLYVAL}/{дайджесты ячеек RELAY, теги аутентификации CGO (гл.~10)},
    2/{Согласование ключей}/{X25519 (Diffie-Hellman)}/{ntor, ntor-v3 (гл.~9)},
    3/Аутентификация/{Ed25519, RSA (постепенно выводится)}/{сертификаты, консенсус (гл.~7, 12)},
    4/Вывод ключей/{HKDF-SHA256, SHAKE-256}/{ключи цепи (гл.~9)}} {
    \node[compito] (c\i) at (0,-\i*1.0) {\c};
    \node[prim] (p\i) at (4.4,-\i*1.0) {\p};
    \draw[->,torgrigio] (c\i) -- (p\i);
    \node[dove] at (6.8,-\i*1.0) {\d};
  }
\end{tikzpicture}

\caption{Пять криптографических задач в Tor, примитивы, которые их решают, и главы, где мы увидим их применение.}
\label{fig:compiti}
\end{figure}

Здравое правило проектирования, которому Tor со временем следует всё строже, --- никогда не использовать один и тот же ключ для разных задач или контекстов.

% ---------------------------------------------------------------------------
\section{Шифры}\label{sec:cripto-cifrari}

\subsection{AES в режиме счётчика}

Почти всё в Tor шифруется с помощью \textbf{AES}, стандартного блочного шифра с блоками по 16 байт и ключами длиной 128 или 256 бит. Сам по себе блочный шифр обрабатывает по одному блоку. Для шифрования потоков произвольной длины Tor использует \textbf{режим счётчика} (CTR): AES шифрует значение счётчика, увеличивающегося на единицу для каждого блока, а полученный \emph{поток ключа} объединяется с данными операцией XOR:
\[ C = P \oplus \operatorname{Trunc}_{|P|}\bigl(\mathrm{AES}_K(\mathit{IV}) \,\|\, \mathrm{AES}_K(\mathit{IV}+1) \,\|\, \cdots\bigr) \]
Расшифрование выполняется той же операцией. Два свойства режима счётчика необходимо помнить всегда:
\begin{itemize}
  \item одна и та же пара ключ--счётчик \textbf{никогда не должна использоваться повторно}: два сообщения, зашифрованные одним потоком ключа, раскрывают XOR своих открытых текстов;
  \item режим счётчика \textbf{пластичен}: изменение бита шифротекста изменяет соответствующий бит открытого текста без знания ключа.
\end{itemize}
AES-128-CTR шифрует RELAY-ячейки в традиционной схеме (\cref{sec:tor1}); его пластичность делает возможными атаки маркирования (tagging). AES-256-CTR шифрует сообщения ntor-v3 (\cref{sec:ntor3}) и дескрипторы onion-сервисов (\cref{sec:os-descrittore}). AES-128 является базовым компонентом CGO (\cref{sec:cgo}).

\subsection{TLS}

Канал между двумя узлами защищён TLS (\cref{cap:07-canali}), который независимо согласует аутентифицированные шифры, например AES-GCM или ChaCha20-Poly1305. Однако Tor не полагается на TLS для подтверждения идентичности реле: реле доказывают её с помощью собственных сертификатов Tor.

% ---------------------------------------------------------------------------
\section{Хеш-функции}\label{sec:cripto-hash}

Криптографическая хеш-функция преобразует вход произвольной длины в отпечаток фиксированной длины. Она проектируется так, чтобы вычислительно невозможно было найти два входа с одинаковым отпечатком (\emph{устойчивость к коллизиям}) либо вход, соответствующий заданному отпечатку. Tor использует четыре семейства, отражающие историю его развития:
\begin{center}\small
\begin{tabularx}{\textwidth}{@{}l c >{\raggedright\arraybackslash}X@{}}
\toprule
Функция & Выход & Где используется \\
\midrule
SHA-1 & 20 байт & RSA-отпечатки реле, накопительный дайджест RELAY-ячеек (tor1), KDF-TOR; считается слабой относительно коллизий, но сохраняется там, где это свойство не требуется \\
SHA-256 & 32 байта & ntor и его KDF, сертификаты Ed25519 и TLS, подписи консенсуса, отпечатки микродескрипторов \\
SHA3-256, SHAKE-256 & 32 байта / переменная длина & ntor-v3, onion-сервисы v3 (адреса, ослеплённые ключи, кольцо HSDir, дескрипторы), дайджесты цепей рандеву \\
Blake2b & переменная длина & доказательство выполнения работы (proof of work) для onion-сервисов (\cref{cap:15-dos}) \\
\bottomrule
\end{tabularx}
\end{center}

SHA-3 --- не улучшенная версия SHA-2, а совершенно иная конструкция (\emph{губка (sponge)}). SHAKE-256 относится к тому же семейству и является функцией с расширяемым выходом: она выдаёт столько байт, сколько требуется, поэтому удобна для выработки ключей. В новых протоколах Tor каждое применение SHA-3 предваряется меткой домена, например \file{"ntor3-curve25519-sha3\_256-1:kdf\_phase1"}. Это не позволяет хешу, вычисленному для одной цели, совпасть с хешем, вычисленным для другой.

\begin{approfondimento}[Накопительные хеши]
Многие библиотеки позволяют приложению добавлять данные в хеш постепенно и в любой момент получать дайджест, не начиная вычисление заново. Tor использует это свойство для \emph{накопительного дайджеста} RELAY-ячеек (\cref{sec:riconoscere}), который обобщает все ячейки, переданные с данным переходом. В Python объект \texttt{hashlib.sha1()} и его методы \texttt{update}, \texttt{digest} и \texttt{copy} делают именно это.
\end{approfondimento}

% ---------------------------------------------------------------------------
\section{MAC и выработка ключей}\label{sec:cripto-mac}

\subsection{HMAC}

MAC --- это отпечаток, который также зависит от ключа: вычислить или проверить его может только тот, кто знает ключ. \textbf{HMAC} строит MAC на основе хеш-функции. Tor использует HMAC-SHA256 в протоколе согласования ключей ntor (\cref{sec:ntor}) и в его функции выработки ключей. В протоколах на основе SHA-3 MAC можно получить проще: хешируя ключ, которому предшествует его длина, а затем сообщение. Длина однозначно указывает, где заканчивается ключ и начинается сообщение~\spec{rend-spec}{Cryptographic building blocks}. В ntor-v3 длина также ставится перед меткой домена (\cref{sec:ntor3}).

Практическая деталь: MAC всегда следует сравнивать за \emph{постоянное время} (в Python --- с помощью \file{hmac.compare\_digest}); иначе время сравнения способно раскрыть, сколько байт совпало.

\subsection{Функции выработки ключей}

Из общего секрета требуется получить несколько ключей: один для шифрования в прямом направлении, один для обратного и два для контроля целостности (\cref{sec:circ-chiavi}). KDF расширяет секрет в поток необходимой длины, из которого извлекаются отдельные ключи. Tor использует три варианта:
\begin{description}
  \item[KDF-TOR] $\mathrm{SHA1}(K_0 \| 0) \,\|\, \mathrm{SHA1}(K_0 \| 1) \,\|\, \ldots$, историческая конструкция, до сих пор используемая \cella{CREATE\_FAST};
  \item[HKDF-SHA256] стандартная конструкция RFC~5869, используемая ntor;
  \item[SHAKE-256] используется напрямую как KDF в ntor-v3 и onion-сервисах.
\end{description}

% ---------------------------------------------------------------------------
\section{Обмен ключами: X25519}\label{sec:cripto-dh}

Обмен ключами Диффи--Хеллмана позволяет двум сторонам установить общий секрет, передавая друг другу только открытые значения. В Tor используется эллиптическая версия \textbf{X25519} (обмен ключами на \emph{Curve25519}). У каждой стороны есть 32-байтовый закрытый ключ $x$ и открытый ключ $X = x \cdot G$, где $G$ --- фиксированная точка кривой, а произведение обозначает умножение точки на целое число. Имея пары $x, X$ и $y, Y$, обе стороны вычисляют одно и то же значение
\[ x \cdot Y = y \cdot X = xy \cdot G, \]
которое наблюдатель, знающий только $X$ и $Y$, получить не может. Ключи имеют длину 32 байта, а обмен выполняется очень быстро. Эти два преимущества привели Tor к отказу от 1024-битного модульного Диффи--Хеллмана, применявшегося в протоколе TAP (\cref{sec:tap}).

Ключ, используемый только для одного обмена и затем удаляемый, называется \textbf{эфемерным}. Именно сочетание эфемерных и долговременных ключей позволяет протоколам согласования ключей Tor обеспечивать одновременно аутентификацию и \emph{прямую секретность} (\cref{sec:circ-requisiti}).

% ---------------------------------------------------------------------------
\section{Подписи: Ed25519 и RSA}\label{sec:cripto-firme}

Цифровая подпись создаётся закрытым ключом и проверяется соответствующим открытым ключом. Tor в основном использует \textbf{Ed25519} --- схему подписи на кривой, эквивалентной Curve25519, с 32-байтовыми открытыми ключами и 64-байтовыми подписями: для идентичностей реле, сертификатов Tor (\cref{sec:canale-certs}) и дескрипторов onion-сервисов. У Ed25519 есть свойство, которое Tor использует необычным способом: умножение открытого ключа $A$ на коэффициент $h$ даёт другой допустимый открытый ключ $A' = h \cdot A$, соответствующий закрытый ключ которого может вычислить тот, кто знает исходный закрытый ключ. На этом основаны ослеплённые ключи onion-сервисов (\cref{sec:os-accecate}).

\textbf{RSA} с 1024-битными ключами --- наследие ранних лет Tor: это устаревшие RSA-идентичности реле и ключи, которыми доверенные серверы каталогов (directory authorities) подписывают консенсус. Ключи серверов каталогов длиннее, и для двух типов действуют разные требования: спецификация требует не менее \textbf{3072} бит для ключа идентичности (\file{KP\_auth\_id\_rsa}) и не менее 2048 для ключа подписи (\file{KP\_auth\_sign\_rsa})~\spec{dir-spec}{Directory authority operation and formats}. Tor постепенно выводит RSA из употребления.

% ---------------------------------------------------------------------------
\section{Современная криптография}\label{sec:cripto-recente}

В последних разработках появляются ещё две специализированные конструкции:
\begin{itemize}
  \item \textbf{POLYVAL}, универсальная хеш-функция, также используемая в AES-GCM-SIV, вместе с AES образует один из строительных блоков CGO --- новой схемы шифрования RELAY-ячеек (\cref{sec:cgo});
  \item \textbf{Elligator2} кодирует точку Curve25519 так, чтобы она была неотличима от случайной строки; obfs4 использует его для сокрытия открытых ключей обмена (\cref{sec:cens-strategie}).
\end{itemize}
\emph{Постквантовая} криптография, которой придётся заменить X25519, когда квантовые компьютеры станут практической угрозой, является активным направлением исследований Tor и рассматривается в \cref{cap:31-futuro}.

% ---------------------------------------------------------------------------
\section*{Итоги}
\begin{itemize}
  \item AES в режиме счётчика шифрует данные; один поток ключа нельзя использовать повторно, а сам режим является пластичным.
  \item SHA-1, SHA-256, SHA3-256/SHAKE-256 и Blake2b сосуществуют в Tor: выбор зависит от времени создания протокола и назначения хеша.
  \item HMAC обеспечивает аутентификацию сообщений, а KDF (KDF-TOR, HKDF, SHAKE) служат для выработки ключей; MAC необходимо сравнивать за постоянное время.
  \item X25519 используется для обмена ключами во всех современных протоколах согласования ключей; Ed25519 подписывает идентичности, сертификаты и дескрипторы; RSA постепенно выводится из употребления.
\end{itemize}

% ---------------------------------------------------------------------------
\section*{Упражнения}

\begin{esercizio}[1]\label{es:05-xor}
При AES-CTR противник знает открытый текст $P_1$ одного сообщения и перехватывает его шифротекст $C_1$, а также шифротекст $C_2$ другого сообщения, ошибочно зашифрованного тем же потоком ключа. Как противник может восстановить $P_2$?
\end{esercizio}

\begin{esercizio}[1]\label{es:05-shake}
В Python вычислите \file{hashlib.shake\_256(b"Tor").hexdigest(16)} и \texttt{hexdigest(32)}. Как связаны два результата? Что это означает для использования SHAKE в качестве KDF?
\end{esercizio}

\begin{esercizio}[2]\label{es:05-dh}
С помощью библиотеки \texttt{cryptography} сгенерируйте две пары ключей X25519 и убедитесь, что они дают одинаковые общие секреты. Что произойдёт, если одна сторона использует 32 нулевых байта как открытый ключ другой стороны? Свяжите результат с проверкой ntor на \enquote{точку на бесконечности}.
\end{esercizio}

\begin{esercizio}[2]\label{es:05-etichette}
Почему протоколы Tor на основе SHA-3 ставят перед каждым входом метку домена? Опишите возможную атаку, если два разных применения одной функции не были бы разделены.
\end{esercizio}

% ---------------------------------------------------------------------------
\section*{Дополнительная литература}
\begin{itemize}
  \item J.-P.~Aumasson, \titolo{Serious Cryptography} (No Starch Press): современное практическое введение во все примитивы этой главы.
  \item D.~J.~Bernstein et al., исходные статьи о Curve25519 и Ed25519, а также RFC~7748 (X25519) и 8032 (Ed25519).
  \item Раздел \enquote{Preliminaries} спецификации Tor~\cite{torspec}, содержащий соглашения о шифрах и хеш-функциях.
\end{itemize}
